\documentclass[a4paper, 12pt]{article}
\usepackage[utf8]{inputenc}
\usepackage[russian]{babel}
\usepackage[T2A]{fontenc}
\usepackage[left=10mm, top=20mm, right=10mm, bottom=15mm]{geometry}
\usepackage{indentfirst}
\usepackage{amsmath}
\usepackage{graphicx}
\usepackage{float}
\usepackage{caption}
\usepackage{array}
\graphicspath{ {graphic/} }
\captionsetup[figure]{name=Рисунок}
  
\title{Отчет по лабораторной работе 2.2.6 \\ Определение энергии активации по температурной зависимости вязкости жидкости}
\author{Максим Осипов, Б03-504}
\date{\today}

\begin{document}
\maketitle

\section*{Аннотация}
\textbf{Цель работы:} 1) измерение скорости падения шариков в жидкости при различных температурах; 2) вычисление вязкости жидкости по закону Стокса; 3) Расчет энергии активации.

\textbf{Оборудование:} стеклянный цилиндр с исследуемой жидкостью, термостат, секундомер, горизонтальный компаратор или микроскоп для измерения диаметра шариков, набор стальных или стеклянных шариков, термометр.

\section{Теоретическая справка}

\subsection{Природа вязкости жидкостей}
В отличие от газов, где вязкость обусловлена переносом импульса при столкновениях молекул, в жидкостях механизм вязкости связан с преодолением молекулами потенциальных барьеров, создаваемых соседними молекулами. Молекула в жидкости колеблется около положения равновесия в потенциальной яме. Для перехода в соседнее положение ("дырку") молекула должна получить дополнительную энергию, достаточную для преодоления барьера. Эта энергия называется \textbf{энергией активации} $W$.

Вероятность такого перехода пропорциональна фактору Больцмана $\exp(-W/kT)$. Поскольку вязкость $\eta$ обратно пропорциональна частоте таких переходов, её температурная зависимость описывается законом Аррениуса-Френкеля:
\[
\eta = A \exp\left(\frac{W}{kT}\right), \qquad (1)
\]
где $A$ — константа, зависящая от свойств жидкости, $k$ — постоянная Больцмана, $T$ — абсолютная температура.
Логарифмируя (1), получаем линейную зависимость $\ln \eta$ от $1/T$:
\[
\ln \eta = \ln A + \frac{W}{k} \cdot \frac{1}{T}. \qquad (2)
\]
По угловому коэффициенту этой прямой можно определить энергию активации $W$.

\subsection{Метод Стокса}
Для определения вязкости используется метод падающего шарика (метод Стокса). На шарик радиусом $r$, движущийся в вязкой жидкости со скоростью $v$, действует сила сопротивления, которая при ламинарном обтекании (малые числа Рейнольдса) определяется законом Стокса:
\[
F_{\text{сопр}} = 6\pi \eta r v. \qquad (3)
\]

При установившемся движении (равномерном падении) сила сопротивления уравновешивается силой тяжести за вычетом силы Архимеда:
\[
6\pi \eta r v_{\text{уст}} = mg - F_A = \frac{4}{3}\pi r^3 \rho g - \frac{4}{3}\pi r^3 \rho_{\text{ж}} g,
\]
где $\rho$ — плотность материала шарика, $\rho_{\text{ж}}$ — плотность жидкости, $g$ — ускорение свободного падения.

Решая это уравнение относительно вязкости, получаем рабочую формулу:
\[
\eta = \frac{2}{9} g r^2 \frac{(\rho - \rho_{\text{ж}})}{v_{\text{уст}}}. \qquad (4)
\]
Установившаяся скорость $v_{\text{уст}} = \frac{l}{t}$ определяется по времени падения $t$ шарика между двумя метками на расстоянии $l$.

\subsection{Число Рейнольдса}
Применимость формулы Стокса ограничена условием ламинарного обтекания, которое характеризуется малым значением числа Рейнольдса для шарика:
\[
Re = \frac{2r v_{\text{уст}} \rho_{\text{ж}}}{\eta} \ll 1 \quad (\text{обычно } Re < 0.5). \qquad (5)
\]

\subsection{Время релаксации}
Начальный участок движения шарика до достижения постоянной скорости характеризуется временем релаксации $\tau$, которое можно оценить по формуле:
\[
\tau = \frac{2}{9} \frac{r^2 \rho}{\eta}. \qquad (6)
\]
Расстояние, проходимое шариком за это время, $S \approx v_{\text{уст}} \tau$. Верхняя метка должна находиться ниже уровня жидкости на расстоянии, значительно превышающем $S$, чтобы измерения проводились в области установившейся скорости.

\section{Экспериментальная установка}
Схема установки представлена на рис. \ref{fig:setup}.

\begin{figure}[H]
\centering
\includegraphics[width=0.5\linewidth]{установка.png} % Имя файла рисунка
\caption{Схема установки для определения вязкости жидкости методом Стокса.}
\label{fig:setup}
\end{figure}

Основные элементы:
\begin{enumerate}
    \item \textbf{Стеклянный цилиндрический сосуд (B)} с исследуемой жидкостью (глицерин). На сосуде нанесены две метки на известном расстоянии $l$.
    \item \textbf{Термостатирующая рубашка (D)}. Через неё циркулирует вода из термостата, что позволяет поддерживать и изменять температуру жидкости в широком диапазоне (от комнатной до $60^\circ$C).
    \item \textbf{Термостат} с регулятором температуры и насосом для циркуляции воды.
    \item \textbf{Набор шариков} (стальных или стеклянных) малого диаметра ($\sim 1$ мм).
    \item \textbf{Секундомер} для измерения времени падения шарика между метками.
    \item \textbf{Микроскоп или компаратор} для точного измерения диаметра шариков.
    \item \textbf{Термометр} для контроля температуры жидкости (или используется шкала термостата).
\end{enumerate}
Плотность глицерина $\rho_{\text{ж}}$ для разных температур определяется по градуировочному графику, прилагаемому к установке. Плотность материала шариков $\rho$ берётся из таблиц ($\rho_{\text{ст}} = 7.7$ г/см$^3$, $\rho_{\text{стекла}} = 2.6$ г/см$^3$).

\section{Обработка данных}
В таблице~\ref{tab:results} представлены результаты обработки экспериментальных данных.  
Расстояние между метками \(l = 0{,}1\ \text{м}\).  
Плотность стекла \(\rho_{\text{ст}} = 2500\ \text{кг/м}^3\), стали \(\rho_{\text{ст}} = 7800\ \text{кг/м}^3\).  
Плотность глицерина при разных температурах принята по справочным данным:
\begin{itemize}
    \item \(t = 20^\circ\text{C}\): \(\rho_{\text{ж}} = 1260\ \text{кг/м}^3\)
    \item \(t = 25^\circ\text{C}\): \(\rho_{\text{ж}} = 1255\ \text{кг/м}^3\)
    \item \(t = 30^\circ\text{C}\): \(\rho_{\text{ж}} = 1250\ \text{кг/м}^3\)
    \item \(t = 35^\circ\text{C}\): \(\rho_{\text{ж}} = 1245\ \text{кг/м}^3\)
    \item \(t = 50^\circ\text{C}\): \(\rho_{\text{ж}} = 1230\ \text{кг/м}^3\)
\end{itemize}

\subsection*{1. Скорость установившегося падения шарика}
\[
v = \frac{l}{t},
\]
где \(l = 0{,}1\ \text{м}\) – расстояние между метками, \(t\) – время падения.

\textbf{Пример:}
\[
v = \frac{0{,}1}{47{,}91} = 0{,}002087\ \text{м/с}.
\]

\subsection*{2. Вязкость жидкости (закон Стокса)}
\[
\eta = \frac{2}{9}\, g\, r^{2}\, \frac{\rho - \rho_{\text{ж}}}{v},
\]
где \(g = 9{,}81\ \text{м/с}^{2}\), \(\rho\) – плотность материала шарика (\(\rho_{\text{ст}} = 2500\ \text{кг/м}^{3}\) для стекла).

\textbf{Пример:}
\[
\eta = \frac{2}{9} \cdot 9{,}81 \cdot (0{,}001000)^{2} \cdot \frac{2500 - 1260}{0{,}002087} = 1{,}295\ \text{Па·с}.
\]

\subsection*{3. Число Рейнольдса (для шарика радиусом \(r\))}
\[
Re = \frac{\rho_{\text{ж}} \, v \, r}{\eta},
\]
где \(\rho_{\text{ж}}\) – плотность жидкости, \(\eta\) – вязкость (полученная в том же опыте).

\textbf{Пример:}
\[
Re = \frac{1260 \cdot 0{,}002087 \cdot 0{,}001000}{1{,}295} = 0{,}00203.
\]

\subsection*{4. Время релаксации}
Время, за которое скорость шарика достигает практически установившегося значения, определяется по формуле:
\[
\tau = \frac{2}{9} \frac{r^{2} \rho}{\eta},
\]
где \(r\) – радиус шарика, \(\rho\) – плотность материала шарика, \(\eta\) – вязкость жидкости.

\textbf{Пример (измерение №1, стекло):}
\[
\tau = \frac{2}{9} \cdot \frac{(0{,}001000)^{2} \cdot 2500}{1{,}295} = 4{,}29 \times 10^{-4}\ \text{с}.
\]

\subsection*{5. Путь релаксации}
Расстояние, проходимое шариком за время релаксации, приближённо равно:
\[
S = v_{\text{уст}} \cdot \tau,
\]
где \(v_{\text{уст}}\) – установившаяся скорость падения.

\textbf{Пример (измерение №1):}
\[
S = 0{,}002087 \cdot 4{,}29 \times 10^{-4} = 8{,}95 \times 10^{-7}\ \text{м} \approx 0{,}9\ \text{мкм}.
\]

\textit{Полученное значение \(S\) пренебрежимо мало по сравнению с расстоянием от поверхности жидкости до верхней метки (\(\approx 0{,}2\)–\(0{,}3\) м), что подтверждает: к моменту прохождения первой метки движение шарика было равномерным.}



\begin{table}[H]
\centering
\caption{Результаты расчётов скорости, вязкости и числа Рейнольдса}
\label{tab:results}
\small % уменьшаем шрифт для надёжности
\begin{tabular}{|c|c|c|c|c|c|c|c|c|}
\hline
№ & Материал & \(d\), мм & \(r\), м & \(t\), с & \(v\), м/с & \(\eta\), Па·с & \(Re\) & \(T\), °C \\
\hline
1  & стекло & 2,00 & 0,001000 & 47,91 & 0,002087 & 1,295 & 0,00203 & 20 \\
2  & стекло & 1,90 & 0,000950 & 47,12 & 0,002122 & 1,150 & 0,00221 & 20 \\
3  & сталь  & 0,78 & 0,000390 & 52,64 & 0,001900 & 1,142 & 0,000818 & 20 \\
4  & сталь  & 0,79 & 0,000395 & 53,40 & 0,001873 & 1,187 & 0,000785 & 20 \\
5  & стекло & 1,98 & 0,000990 & 33,04 & 0,003027 & 0,879 & 0,00428 & 25 \\
6  & стекло & 2,02 & 0,001010 & 32,94 & 0,003036 & 0,912 & 0,00422 & 25 \\
7  & сталь  & 0,72 & 0,000360 & 37,93 & 0,002637 & 0,701 & 0,00170 & 25 \\
8  & сталь  & 0,60 & 0,000300 & 54,85 & 0,001823 & 0,704 & 0,000975 & 25 \\
9  & стекло & 1,92 & 0,000960 & 22,74 & 0,004397 & 0,571 & 0,00924 & 30 \\
10 & стекло & 1,91 & 0,000955 & 22,72 & 0,004401 & 0,565 & 0,00930 & 30 \\
11 & сталь  & 0,62 & 0,000310 & 30,93 & 0,003233 & 0,424 & 0,00296 & 30 \\
12 & сталь  & 0,82 & 0,000410 & 22,38 & 0,004468 & 0,537 & 0,00426 & 30 \\
13 & стекло & 2,11 & 0,001055 & 15,44 & 0,006477 & 0,470 & 0,0181  & 35 \\
14 & стекло & 2,10 & 0,001050 & 15,69 & 0,006373 & 0,473 & 0,0176  & 35 \\
15 & сталь  & 0,77 & 0,000385 & 21,13 & 0,004733 & 0,447 & 0,00507 & 35 \\
16 & сталь  & 0,85 & 0,000425 & 14,90 & 0,006711 & 0,385 & 0,00922 & 35 \\
17 & стекло & 2,00 & 0,001000 & 7,04  & 0,01420  & 0,195 & 0,0896  & 50 \\
18 & стекло & 2,05 & 0,001025 & 6,92  & 0,01445  & 0,201 & 0,0907  & 50 \\
19 & сталь  & 0,75 & 0,000375 & 9,08  & 0,01101  & 0,183 & 0,0277  & 50 \\
20 & сталь  & 0,90 & 0,000450 & 6,13  & 0,01631  & 0,178 & 0,0507  & 50 \\
\hline
\end{tabular}
\end{table}

Все полученные числа Рейнольдса значительно меньше \(0{,}5\) (максимальное \(\approx 0{,}09\)), что подтверждает ламинарный режим обтекания и корректность применения формулы Стокса.

\subsection{Оценка погрешности измерений}

Инструментальные погрешности:  
\(\Delta d = 0{,}01\ \text{мм}\) (микрометр), \(\Delta t = 0{,}1\ \text{с}\) (секундомер), \(\Delta l = 1\ \text{мм}\) (линейка).  
Относительная погрешность вязкости для типичного измерения (например, №1) составляет:

\[
\frac{\Delta\eta}{\eta} \approx \sqrt{\left(2\frac{\Delta d}{d}\right)^2 + \left(\frac{\Delta t}{t}\right)^2 + \left(\frac{\Delta l}{l}\right)^2 } \approx \sqrt{(0{,}01)^2 + (0{,}0021)^2 + (0{,}01)^2} \approx 0{,}014\ (1{,}4\%).
\]

Для разных измерений погрешность варьируется в пределах \(1\)–\(2\%\). Плотности материалов и глицерина считаются известными с достаточной точностью.

\section{Вычисления}
\subsection*{График зависимости \(\ln \eta\) от \(1/T\)}
По экспериментальным данным (таблица) был построен график зависимости логарифма вязкости от обратной абсолютной температуры (рис. \ref{fig:lneta}). Точками показаны результаты отдельных измерений, сплошная линия — линейная аппроксимация методом наименьших квадратов.

\begin{figure}[h!]
    \centering
    \includegraphics[width=0.8\textwidth]{график.png} % замените на имя вашего файла
    \caption{Зависимость \(\ln \eta\) от \(1/T\).}
    \label{fig:lneta}
\end{figure}

Аппроксимирующая прямая описывается уравнением
\[
\ln \eta = a + b \cdot \frac{1}{T},
\]
где угловой коэффициент \(b = 5700\ \text{К}\), свободный член \(a = -12{,}5\) (значение получено из регрессии). 

\subsection*{Определение энергии активации}
Согласно формуле (2.16):
\[
W = k\,\frac{d(\ln \eta)}{d(1/T)},
\]
где \(k = 1{,}38\times 10^{-23}\) Дж/К — постоянная Больцмана, а производная \(\frac{d(\ln \eta)}{d(1/T)}\) равна угловому коэффициенту \(b\) построенной прямой.

Подставляя найденное значение \(b = 5700\ \text{К}\), получаем энергию активации для одной молекулы:
\[
W = k b = 1{,}38\times 10^{-23} \cdot 5700 = 7{,}87\times 10^{-20}\ \text{Дж}.
\]


Для моля вещества энергия активации выражается через универсальную газовую постоянную \(R = k N_A = 8{,}31\ \text{Дж/(моль·К)}\):
\[
W_{\text{мол}} = R b = 8{,}31 \cdot 5700 = 47\,400\ \text{Дж/моль} = 47{,}4\ \text{кДж/моль}.
\]

Полученное значение хорошо согласуется с типичными энергиями активации вязкого течения для органических жидкостей (20–50 кДж/моль), что подтверждает корректность выполненных измерений и расчётов.

\subsection*{Оценка погрешности определения энергии активации}

Для оценки погрешности энергии активации воспользуемся методом линейной регрессии по средним значениям \(\ln\eta\) при каждой температуре. Исходные данные для пяти температурных точек сведены в таблицу~\ref{tab:means}.

\begin{table}[h!]
\centering
\caption{Средние значения логарифма вязкости и обратной температуры}
\label{tab:means}
\begin{tabular}{|c|c|c|c|}
\hline
\(T,\ K\) & \(1/T,\ 10^{-3}\ \text{К}^{-1}\) & \(\ln\eta\) & \(\sigma_{\ln\eta}\) \\
\hline
293,15 & 3,411 & 0,1757 & 0,0290 \\
298,15 & 3,354 & –0,2319 & 0,0705 \\
303,15 & 3,299 & –0,6528 & 0,0697 \\
308,15 & 3,245 & –0,8158 & 0,0479 \\
323,15 & 3,095 & –1,6655 & 0,0281 \\
\hline
\end{tabular}
\end{table}

По этим точкам методом наименьших квадратов получено уравнение прямой
\[
\ln\eta = a + b\cdot\frac{1}{T},\quad b = 5700\ \text{К}.
\]
Стандартную ошибку углового коэффициента \(\sigma_b\) вычисляем по формуле
\[
\sigma_b = \sqrt{\frac{1}{n-2}\cdot\frac{\sum_{i=1}^{n}(y_i-a-bx_i)^2}{\sum_{i=1}^{n}(x_i-\bar{x})^2}},
\]
где \(n=5\), \(x_i=1/T_i\), \(y_i=\ln\eta_i\). Подставляя значения, получаем
\[
\sum (x_i-\bar{x})^2 = 5{,}85\times 10^{-8},\quad \sum (y_i-a-bx_i)^2 = 0{,}0210,
\]
\[
\sigma_b = \sqrt{\frac{0{,}0210}{3\cdot5{,}85\times10^{-8}}} \approx 346\ \text{К}.
\]

Тогда энергия активации
\[
W = k b = 1{,}38\times10^{-23}\cdot5700 = 7{,}87\times10^{-20}\ \text{Дж}
\]
и её абсолютная погрешность
\[
\Delta W = k\,\sigma_b = 1{,}38\times10^{-23}\cdot346 = 4{,}77\times10^{-21}\ \text{Дж}.
\]


Для одного моля вещества
\[
W_{\text{мол}} = R b = 8{,}31\cdot5700 = 47\,400\ \text{Дж/моль} = 47{,}4\ \text{кДж/моль},
\]
\[
\Delta W_{\text{мол}} = R\,\sigma_b = 8{,}31\cdot346 = 2880\ \text{Дж/моль} \approx 2{,}9\ \text{кДж/моль}.
\]

Таким образом, энергия активации вязкого течения исследованного образца составляет
\[
W_{\text{мол}} = (47{,}4 \pm 2{,}9)\ \text{кДж/моль}.
\]

Полученное значение хорошо согласуется с литературными данными для глицерина (около 50 кДж/моль), а относительная погрешность ~6\% является приемлемой для лабораторного эксперимента.

\section*{Вывод}

В результате выполнения лабораторной работы были проведены измерения скорости падения стальных и стеклянных шариков в глицерине при пяти различных температурах в интервале от 20 до 50 °C. На основе полученных данных с использованием закона Стокса рассчитаны значения вязкости жидкости для каждого измерения. Анализ чисел Рейнольдса (\(Re < 0{,}1\) для всех опытов) подтвердил ламинарный характер обтекания шариков, что обосновывает корректность применения формулы Стокса в условиях эксперимента. Оценка времени релаксации (\(\tau \sim 10^{-4}\) с) и пути релаксации (\(S \sim 1\) мкм) показала, что движение шариков было равномерным уже на начальном участке, поэтому измеренные скорости соответствуют установившимся значениям.

Построена зависимость логарифма вязкости от обратной температуры \(\ln\eta = f(1/T)\), имеющая линейный характер, что согласуется с моделью Аррениуса–Френкеля. Методом наименьших квадратов определён угловой коэффициент прямой \(b = 5700 \pm 346\) К. С помощью соотношения (2.16) вычислена энергия активации вязкого течения глицерина:

* для одной молекулы: \(W = (7{,}87 \pm 0{,}48) \times 10^{-20}\) Дж или \((0{,}49 \pm 0{,}03)\) эВ;
* для одного моля вещества: \(W_{\text{мол}} = (47{,}4 \pm 2{,}9)\) кДж/моль.

Полученное значение энергии активации хорошо согласуется с известными литературными данными для глицерина (около 50 кДж/моль). Относительная погрешность определения величины \(W_{\text{мол}}\) составила около 6\%, что является приемлемым для лабораторного практикума и подтверждает достоверность выполненных измерений и корректность использованной методики обработки.


\end{document}